% ============================================================
% CHAPTER 6: DISCUSSION
% ============================================================
\chapter{Discussion}
\label{ch:discussion}

This chapter interprets the results in relation to the research aims and wider literature. Section~\ref{sec:disc-cyclingforall} evaluates the West Midlands network against LTN 1/20's "cycling for all" principle, situating the connectivity findings within evidence on cycling demand and mode share. Section~\ref{sec:disc-junction} examines junction stress as a structural mechanism of low-stress network severance, distinguishing the barrier roles of crossings and roundabouts. Section~\ref{sec:disc-planning} draws implications from the growth simulation for investment prioritisation and cross-authority coordination. Section~\ref{sec:disc-osvalidation} assesses OS NGD cycle lane data through a first independent comparison with OpenStreetMap, and Section~\ref{sec:disc-limitations} discusses limitations and directions for future work.

\section{Evaluating "Cycling for All": Network Connectivity Against LTN 1/20 Objectives}
\label{sec:disc-cyclingforall}

\section{Junction Stress and the Structural Severance of the Low-Stress Network}
\label{sec:disc-junction}

\section{Network Growth and Implications for Investment Coordination}
\label{sec:disc-planning}

\section{Assessing OS NGD Cycle Lane Data: A First Independent Comparison with OpenStreetMap}
\label{sec:disc-osvalidation}

% TODO: include the honest treatment of the Furth et al. (2026)
% primal-graph junction stress propagation critique; OS NGD
% cycle lane undercount and dataset maturity.

\section{Limitations and Future Work}
\label{sec:disc-limitations}
